\documentclass{article}
\usepackage{booktabs}
\usepackage{pgfplots}
\pgfplotsset{compat=1.18}
\usepackage{geometry}
\geometry{a4paper, margin=1in}

\title{RoRA Experiment Results}
\author{Generated from Experiment Outputs}
\date{\today}

\begin{document}
\maketitle

\section{Experiment 1: MNIST Multi-Task Learning}

\begin{table}[h]
\centering
\caption{MNIST Multi-Task Learning Results}
\label{tab:experiment1}
\begin{tabular}{lcc}
\toprule
Method & Task A (Even/Odd) & Task B (Bit-Parity) \\
\midrule
Base (frozen) & 96.97\% & N/A \\
RoRA Rank 4 & $98.87 \pm 0.05\%$ & $98.78 \pm 0.08\%$ \\
RoRA Rank 8 & $98.91 \pm 0.04\%$ & $98.80 \pm 0.05\%$ \\
RoRA Rank 16 & $98.87 \pm 0.09\%$ & $98.75 \pm 0.13\%$ \\
\bottomrule
\end{tabular}
\end{table}

\begin{figure}[h]
\centering
\begin{tikzpicture}
\begin{axis}[
    width=\textwidth,
    height=0.5\textwidth,
    xlabel={RoRA Rank},
    ylabel={Accuracy (\%)},
    legend pos=south east,
    grid=major,
    xmin=2, xmax=18,
    ymin=95, ymax=100,
    xtick={4, 8, 16},
]
\addplot[mark=*, blue, thick] coordinates {(4,98.87), (8,98.91), (16,98.87)};
\addlegendentry{Task A (Even/Odd)}
\addplot[mark=square*, red, thick] coordinates {(4,98.78), (8,98.80), (16,98.75)};
\addlegendentry{Task B (Bit-Parity)}
\addplot[dashed, gray, thick] coordinates {(2,96.97) (18,96.97)};
\addlegendentry{Base (frozen)}
\end{axis}
\end{tikzpicture}
\caption{Performance comparison for multi-task learning}
\label{fig:experiment1}
\end{figure}

\end{document}